\begin{frame}{왜 알고리즘은 ``1개 이상''인가}
  \begin{itemize}
    \item Block A: 두 번째 알고리즘은 \textbf{갱신 규칙 한 줄의 차이}(6.3) —
      Q-learning vs SARSA.
    \item Block B: 각 알고리즘이 \textbf{별개의 신경망 학습 파이프라인}.
      \begin{itemize}
        \item DQN: 리플레이 버퍼 + 타깃 네트워크가 따라온다.
        \item PPO: GAE + 클리핑이 따라온다.
        \item 모방학습: 시연 데이터셋이 따라온다.
      \end{itemize}
    \item 팀 단위 현실적 예산: ``한 파이프라인을 제대로 돌리고, 여유가
      있으면 비교''.
  \end{itemize}
  \vskip0.6em
  \begin{block}{변하지 않는 것이 핵심}
    시드 $\geq 3$ · 베이스라인(무작위 정책) · 학습/평가 분리는 알고리즘과
    무관한 \textbf{강화학습 실험의 골격}이라 Block A와 동일하다.
  \end{block}
\end{frame}
